\documentclass[11pt]{article}
\usepackage[margin=1in]{geometry}
\usepackage{amsmath,amssymb,amsthm,hyperref}
\newcommand{\R}{\mathbb R}
\newcommand{\C}{\mathbb C}
\newcommand{\dd}{\,\mathrm d}
\newtheorem{theorem}{Theorem}
\title{Independent check of the Polymath heat remainder constant}
\date{20 September 2026}
\begin{document}
\maketitle

\section{Source, scope, and statement}

The source checked here is D.~H.~J. Polymath,
\emph{Effective approximation of heat flow evolution of the Riemann
$\xi$ function, and a new upper bound for the de Bruijn--Newman
constant}, arXiv:1904.12438v2, original author file
\texttt{debruijn.tex} \cite{polymath}. Its SHA-256 is
\texttt{560a28fe31bec92dd793820222e9e73a1fc6958a08344033a946b2ccaba225e5}.
The imported Riemann--Siegel bounds below are the literal statement
of its Proposition \texttt{arias-prop}, source lines 684--715,
and the following explanation and $C_0$ bound through line 724.
This is an independent derivation from those stated bounds, not
a claim to have independently reread or reproved Arias de Reyna's
underlying paper. The complete original proof of
\texttt{RTN-prop}, source lines 743--920, has been read for this check.
The receiving author's appendix \texttt{SOURCE\_REMAINDER\_REPAIR.tex}
was read after the cutoff and Gaussian computations were independently
checked. No failure in its theorem PMH11 was found.

For nonreal $s$, preserve the original function
\begin{equation}\label{defR}
 R_{0,N}(s)=\frac18\frac{s(s-1)}2\pi^{-s/2}\Gamma(s/2)
 \int_{N\swarrow N+1}
 \frac{w^{-s}e^{i\pi w^2}}{e^{\pi iw}-e^{-\pi iw}}\dd w.
\end{equation}
Here $N\swarrow N+1$ has direction $e^{5\pi i/4}$, oriented as
in the source, and crosses between $N$ and $N+1$; the power is
the branch in the source contour formula. For $t>0$ put
\begin{equation}\label{heatR}
 R_{t,N}(s)=\int_\R R_{0,N}(s+\sqrt t\,v)
                      \frac{e^{-v^2}}{\sqrt\pi}\dd v.
\end{equation}
These are source \texttt{RON-def} and the following heat evolution
definition, read together with \texttt{RTN-def} \cite{polymath}.

Use the source's nowhere-zero holomorphic function on
$\C\setminus(-\infty,1]$,
\begin{equation}\label{defM}
 M_0(s)=\frac18\frac{s(s-1)}2\pi^{-s/2}\sqrt{2\pi}
 \exp\!\left((s/2-1/2)\operatorname{Log}(s/2)-s/2\right),
\end{equation}
with its specified holomorphic logarithm and principal logarithm
branches. Its logarithmic derivative is
\begin{equation}\label{alpha}
 \alpha(s)=\frac1{2s}+\frac1{s-1}+\frac12\operatorname{Log}(s/(2\pi)),
 \qquad
 \alpha'(s)=-\frac1{2s^2}-\frac1{(s-1)^2}+\frac1{2s}.
\end{equation}
For $\Im s>3$, the source bound
\begin{equation}\label{alphabound}
 |\alpha'(s)|\le\frac1{2(\Im s-3)}
\end{equation}
follows by bounding the three terms of \eqref{alpha} and summing
the positive geometric series. These are source \texttt{M-def},
\texttt{logM}, \texttt{alpha-form}, and
\texttt{alpha-deriv-bound} \cite{polymath}.

\begin{theorem}\label{main}
Let $0\le\sigma\le1$, $T\ge100$, and $0\le t\le1/2$. Set
\begin{gather*}
 T'=T+\pi t/8,\qquad a=\sqrt{T'/(2\pi)},\qquad N=\lfloor a\rfloor,
 \qquad p=1-2(a-N),\\
 U=\exp\!\left[-i\left(\frac{T'}2\log\frac{T'}{2\pi}
                   -\frac{T'}2-\frac\pi8\right)\right],\\
 C_0(p)=\frac{e^{\pi i(p^2/2+3/8)}-i\sqrt2\cos(\pi p/2)}{2\cos(\pi p)},
\end{gather*}
with removable values at $p=\pm1/2$. Then
\begin{align}\label{result}
 &\left|\frac{R_{t,N}(\sigma+iT)}
 {(-1)^{N-1}Ue^{\pi i\sigma/4}e^{t\pi^2/64}M_0(iT')}-C_0(p)\right|\\
 &\hspace{12mm}\le
 \left(\frac{0.397\,9^\sigma}{a-0.865}+\frac5{3(T-6)}\right)
                      e^{3.49/(T-4)}.\nonumber
\end{align}
For $t>0$ this is precisely the bound stated in source
\texttt{RTN-prop}. The endpoint $t=0$ is established separately below.
\end{theorem}

The proof repairs four source passages: the Stirling numerator $1/3$
in lines 769--776; inconsistent denominators in Gaussian calculations
including lines 792 and 853--857; omission of one possible final
Riemann--Siegel remainder after the choice of $K$ at line 808; and the
changed lower summation cutoff in the equality at line 903. The
last change is allowed as an inequality, but not as that equality.
None of these repairs changes \eqref{result}.

\section{The exact imported coefficient bounds}

For real $u$, $T'>0$, positive integer $K$, and $a,N,p,U$ as above,
Polymath's Proposition \texttt{arias-prop} gives
\begin{equation}\label{RSexp}
 \int_{N\swarrow N+1}\frac{w^{-u-iT'}e^{i\pi w^2}}
 {e^{\pi iw}-e^{-\pi iw}}\dd w
 =(-1)^{N-1}Ua^{-u}
 \left(C_0(p)+\sum_{k=1}^K\frac{C_k(p,u)}{a^k}+RS_K(u+iT')\right).
\end{equation}
All gamma factors and sign domains in the following input are retained:
\begin{align}
 |C_k(p,u)|&\le\frac{\sqrt2}{2\pi}
                      \frac{9^u\Gamma(k/2)}{2^k} &&(u>0),\label{ckpos}\\
 |C_k(p,u)|&\le\frac{2^{1/2-u}}{2\pi}
       \frac{\Gamma(k/2)}{2\pi((3-2\log2)\pi)^{k/2}} &&(u\le0),\label{ckneg}\\
 |RS_K(u+iT')|&\le\frac17 2^{3u/2}
                         \Gamma((K+1)/2)(1.1/a)^{K+1} &&(u\ge0),\label{rspos}\\
 |RS_K(u+iT')|&\le\frac12(9/10)^{\lceil-u\rceil}
                   \Gamma((K+1)/2)(1.1/a)^{K+1}
                          &&(u<0,\ K+u\ge2).\label{rsneg}
\end{align}
The source labels are \texttt{ck-bound-1}, \texttt{ck-bound-2},
\texttt{rsk-bound-1}, \texttt{rsk-bound-2}. Polymath attributes these
to Arias de Reyna, Theorems 3.1, 4.1, 4.2 and equations (3.2), (5.2)
\cite{arias}. The further source input, label \texttt{cop}, is
$|C_0(p)|\le1/2$ for $p\in[-1,1]$, attributed there to the $n=0$
case of Arias de Reyna's Theorem 6.1 \cite{polymath,arias}.

\section{Cutoff admissibility and the previously missing remainder}

Write $b=2a^2=T'/\pi$. We have $b\ge100/\pi>31$ and $a>3$.
Choose the source cutoff
\[
 K_u=1\quad(u\ge0),\qquad
 K_u=\max(\lfloor-u\rfloor+3,\lfloor b\rfloor)\quad(u<0).
\]
This is a measurable, integer-valued function. For $u<0$,
$\lfloor-u\rfloor+u>-1$, hence $K_u+u>2$, which verifies
the only additional condition in \eqref{rsneg}. Define
\[
 S(u)=\sum_{k=1}^{K_u}|C_k(p,u)|a^{-k}+|RS_{K_u}(u+iT')|.
\]
For $u>0$, \eqref{ckpos} and \eqref{rspos} give
\[
 S(u)\le\frac{0.2\,9^u}{a}+\frac{0.173\,2^{3u/2}}{a^2}
 \le0.2\,9^u\left(\frac1a+\frac{0.865}{a^2}\right)
 \le\frac{0.2\,9^u}{a-0.865}.
\]
Indeed $\sqrt{2\pi}/(4\pi)<0.2$, $1.21/7<0.173$,
$2^{3/2}<9$, and $0.2(0.865)=0.173$.
At $u=0$, use \eqref{ckneg}; its $k=1$ coefficient is below
$0.036(0.445)\sqrt\pi<0.0285<0.2$. Thus the same bound holds
also at $u=0$, without extending the domain of \eqref{ckpos}.

For $u<0$, put $q=0.445$. The exact constants in \eqref{ckneg}
satisfy
\[
 \frac{\sqrt2}{4\pi^2}<0.036,\qquad
 ((3-2\log2)\pi)^{-1/2}<q.
\]
The identity
\[
 \frac{\Gamma((k+2)/2)}{a^{k+2}}
 =\frac{k}{b}\frac{\Gamma(k/2)}{a^k}
\]
shows that on each parity subsequence, as long as both indices remain
at most $b$, successive terms $q^k\Gamma(k/2)/a^k$ have ratio at
most $q^2$. Consequently
\begin{align*}
 \sum_{1\le k\le b}\frac{|C_k(p,u)|}{a^k}
 &\le0.036\,2^{-u}\left(
 \frac{q\sqrt\pi}{a}+
 \frac{q^2}{1-q^2}\frac1{a^2}+
 \frac{q^3\sqrt\pi}{2(1-q^2)}\frac1{a^3}\right)\\
 &\le0.036\,2^{-u}\left(
 \frac{0.445\sqrt\pi}{a}+\frac{0.247}{a^2}+\frac{0.098}{a^3}\right)
 \le\frac{0.0285\,2^{-u}}{a-0.353}.\tag{L}\label{low}
\end{align*}
For the last step, expand $(a-0.353)^{-1}$ as a positive geometric
series and compare its first three terms, using
\[
 0.036(0.445)\sqrt\pi<0.0285,\quad
 0.036(0.247)<0.0285(0.353),\quad
 0.036(0.098)<0.0285(0.353)^2.
\]

If $K_u=\lfloor b\rfloor$, its remainder index is
$r=K_u+1=\lfloor b\rfloor+1\ge32$. Since
$a^2\ge(r-1)/2$, \eqref{rsneg} gives
\begin{equation}\label{vr}
 a2^u|RS_{K_u}(u+iT')|
 \le v_r:=\frac12(1.1)^r\Gamma(r/2)
                           ((r-1)/2)^{(1-r)/2}.
\end{equation}
Here $2^u\le1$ and $(9/10)^{\lceil-u\rceil}\le1$.
Direct gamma recurrence gives the exact ratio
\[
 \frac{v_{r+2}}{v_r}=1.21\frac r{r+1}
     \left(\frac{r-1}{r+1}\right)^{(r-1)/2}
 \le1.21\exp\!\left(-\frac{r-1}{r+1}\right)
 \le1.21e^{-31/33}<1.
\]
The first inequality uses $\log(1-x)\le-x$ with $x=2/(r+1)$;
the last follows already from $e^{31/33}>1+31/33>1.21$.
The exact squared base expressions are
\[
 v_{32}^2=\frac14(11/10)^{64}(15!)^2(2/31)^{31},\qquad
 v_{33}^2=\frac14(11/10)^{66}
 \left(\frac{32!}{4^{16}16!}\right)^2\pi\,16^{-32}.
\]
Their certified upper bounds are respectively
$(4.897\cdot10^{-6})^2$ and $(3.268\cdot10^{-6})^2$.
Both are less than $(0.0005)^2$. The two parity inductions now prove
\[
 |RS_{K_u}(u+iT')|\le\frac{0.0005\,2^{-u}}a
 \le\frac{0.0005\,2^{-u}}{a-0.353}
 \quad\hbox{when }K_u=\lfloor b\rfloor.
\]
This is the remainder absent from the source's displayed bound
after \texttt{laf}. Together with \eqref{low}, it fits
$0.029\,2^{-u}/(a-0.353)$.

If $K_u>\lfloor b\rfloor$, necessarily $K_u=\lfloor-u\rfloor+3$.
Every coefficient index above $b$ satisfies $b\le k\le-u+4$;
the final remainder index $K_u+1$ also satisfies those inequalities.
By \eqref{ckneg}, its coefficient majorant is at most
$(1/2)2^{-u}(1.1)^k\Gamma(k/2)/a^k$, and \eqref{rsneg}
puts the remainder in the same sum. The coefficient indices and
the remainder index are distinct, so no index is counted twice
for this upper bound. If $b$ is an integer, the extra index $k=b$
in the enlarged high sum is harmless. Ties
$\lfloor-u\rfloor+3=\lfloor b\rfloor$ were already covered by
\eqref{vr}. Therefore, for every real $u$,
\begin{equation}\label{Sbound}
 S(u)\le\frac{0.2\,9^u}{a-0.865}
       +\frac{0.029\,2^{-u}}{a-0.353}
       +\frac{2^{-u}}2\sum_{b\le k\le-u+4}
                              \frac{(1.1)^k\Gamma(k/2)}{a^k}.
\end{equation}
All sums here and below are over integers. Empty sums are zero.

\section{The pointwise exponent and the exact Gaussian integral}

Let $D=T-3\ge97$ and $c=11/12$. Source Lemma
\texttt{elem-lem}(v), attributed to Boyd \cite{polymath,boyd},
states an exponential Stirling error at $z=(u+iT')/2$ of modulus
at most $1/(12(|z|-0.33))$. Thus its contribution is at most
$1/(6(T'-0.66))$. Taylor's formula using \eqref{alphabound}
and the horizontal segment from $iT'$ to $u+iT'$ bounds the
logarithmic Taylor remainder by $u^2/(4(T'-3))$. Finally
\[
 \alpha(iT')=\log a+\frac{\pi i}4+\frac1{2iT'}+\frac1{iT'-1},
 \qquad
 \left|\frac1{2iT'}+\frac1{iT'-1}\right|\le\frac3{2T'}.
\]
There is consequently a complex exponential error $E(u)$ with
\begin{equation}\label{Ebound}
 |E(u)|\le\frac{u^2}{4(T'-3)}+\frac{3|u|}{2T'}
                      +\frac1{6(T'-0.66)}
 \le\frac{u^2+6|u|+2/3}{4D}
 \le\frac{u^2+c}{D}.
\end{equation}
The last step uses $6|u|\le3u^2+3$. In particular the Stirling
numerator is $2/3$, not the source's $1/3$.

For $t>0$, the source contour shift \texttt{RTN-def} with
$\beta_N=\pi i/4$ reads
\[
 R_{t,N}(\sigma+iT)=e^{t\pi^2/64}
 \int_\R e^{-\sqrt t\,v\pi i/4}
 R_{0,N}(\sigma+\sqrt t\,v+iT')
                         \frac{e^{-v^2}}{\sqrt\pi}\dd v.
\]
The shift is admissible because $T,T'>0$; the source justifies
this by slower-than-Gaussian growth on the connecting horizontal
strips. Set $u=\sigma+\sqrt t\,v$. Formula \eqref{RSexp} contains
$a^{-u}$, which cancels the factor $e^{u\log a}$ from $M_0$.
The two phases $e^{\pi iu/4}$ and $e^{-\sqrt t\,v\pi i/4}$
leave exactly $e^{\pi i\sigma/4}$. The real probability measure is
\[
 \dd\nu(u)=\frac1{\sqrt{\pi t}}e^{-(u-\sigma)^2/t}\dd u,
 \qquad \int_\R\dd\nu=1.
\]
Using $|C_0|\le1/2$ and $|e^z-1|\le e^{|z|}-1$, the left
side of \eqref{result} is bounded by
\begin{equation}\label{epsdelta}
 \frac12\epsilon+\delta,\qquad
 \epsilon=\int_\R(e^{(u^2+c)/D}-1)\dd\nu(u),\qquad
 \delta=\int_\R e^{(u^2+c)/D}S(u)\dd\nu(u).
\end{equation}

For every real $l$, completing the square gives exactly
\begin{equation}\label{J}
 J(l):=\int_\R e^{(u^2+c)/D+2lu}\dd\nu(u)
 =e^{2l\sigma+tl^2}(1-t/D)^{-1/2}
       \exp\!\left(\frac cD+\frac{(\sigma+tl)^2}{D-t}\right).
\end{equation}
To check the algebra, the exponent before integration is
\[
 -\frac{D-t}{tD}u^2+2(\sigma/t+l)u+c/D-\sigma^2/t.
\]
Its Gaussian integral contributes $(1-t/D)^{-1/2}$ and exponent
$c/D-\sigma^2/t+tD(\sigma/t+l)^2/(D-t)$.
Subtracting $2l\sigma+tl^2+c/D$ leaves
$(\sigma+tl)^2/(D-t)$, proving \eqref{J}. Since $D>t$,
all these integrals converge. No occurrence of $D$ in this identity
can be replaced by $T'-0.33$.

The elementary inequality $-\log(1-x)\le x/(1-x)$ for
$0\le x<1$ yields
\[
 \log J(0)\le\frac{c+\sigma^2+t/2}{D-t}
 \le\frac{13/6}{T-3.5}\le\frac{10}{3(T-6)}.
\]
The last inequality follows from $13(T-6)\le20(T-3.5)$.
Since $e^v-1\le ve^v$ for $v\ge0$ and
\[
 \frac{10}{3(T-6)}\le\frac{3.49}{T-4}
 \quad\Longleftrightarrow\quad 0.47T\ge22.82,
\]
we obtain
\begin{equation}\label{epsilonfinal}
 \frac12\epsilon=\frac12(J(0)-1)
       \le\frac5{3(T-6)}e^{3.49/(T-4)}.
\end{equation}

\section{The first two integrated coefficient terms}

By \eqref{Sbound} write $\delta\le\delta_1+\delta_2+\delta_3$,
where $\delta_1=0.2J(\log3)/(a-0.865)$,
$\delta_2=0.029J(-\log\sqrt2)/(a-0.353)$, and $\delta_3$
integrates the final sum in \eqref{Sbound}.
Set $L=\log3$ and
\[
 A=c+(1+L/2)^2+1/4=3.5670161955\ldots.
\]
After removing $9^\sigma e^{tL^2}$ from \eqref{J}, the remaining
logarithm is bounded by
\[
 \frac cD+\frac{(\sigma+tL)^2+t/2}{D-t}
 \le\frac A{T-3.5}\le\frac A{T-4}.
\]
Since $A>3.49$ and $T-4\ge96$,
\[
 0.2e^{tL^2}e^{(A-3.49)/(T-4)}
 \le0.2\exp\!\left(L^2/2+(A-3.49)/96\right)
 <0.365986<0.366.
\]
The strict numerical comparison is proved by exact rational series
certificates in the last section; it is not a rounded floating-point
comparison. Therefore
\begin{equation}\label{deltaone}
 \delta_1\le\frac{0.366\,9^\sigma}{a-0.865}e^{3.49/(T-4)}.
\end{equation}
This absorption is necessary: merely replacing $5/6$ by $11/12$
in the source's old exponent and claiming it is at most $3.49$
would not be valid.

Let $M=\log\sqrt2$. As $0<M<1$, $0\le\sigma\le1$ and
$0\le t\le1/2$ imply $|\sigma-tM|\le1$. After removing
$2^{-\sigma}e^{tM^2}$ from \eqref{J}, the remaining logarithm
is at most
\[
 \frac cD+\frac{1+t/2}{D-t}
 \le\frac{13/6}{T-3.5}<\frac{3.49}{T-4}.
\]
The rational-series comparison $0.029e^{M^2/2}<0.031$ now gives
\begin{equation}\label{deltatwo}
 \delta_2\le\frac{0.031\,2^{-\sigma}}{a-0.353}e^{3.49/(T-4)}.
\end{equation}

\section{The exact tail cutoff and its enlargement}

Tonelli's theorem for nonnegative measurable functions gives the
exact expression
\begin{equation}\label{tail}
 \delta_3=\frac1{2\sqrt{\pi t}}
 \sum_{k\ge b}\frac{(1.1)^k\Gamma(k/2)}{a^k}
 \int_{-\infty}^{4-k}
 \exp\!\left(\frac{u^2+c}D-\frac{(u-\sigma)^2}t-u\log2\right)\dd u.
\end{equation}
The cutoff is exactly $b=T'/\pi$, not $T'/(2.2\pi)$.
Here $k\ge32$ and $u\le-28$. To prepare an upper bound enlarged
to $k\ge14$, it is enough to use $u\le-10$. Because $\sigma\ge0$,
$(u-\sigma)^2\ge u^2$. Moreover, $u^2\ge100$, $D\ge97$,
and $1/(2t)\ge1$ imply
\[
 \frac{u^2+c}D
 \le u^2\frac{1+c/100}{97}\le\frac{u^2}{2t}.
\]
Thus the exponent in \eqref{tail} is at most
$-u^2/(2t)-u\log2$. For $u\le4-k$,
$-u^2/(2t)\le(k-4)u/(2t)$, so
\begin{align*}
 \int_{-\infty}^{4-k}e^{-u^2/(2t)-u\log2}\dd u
 &\le\frac{e^{-(k-4)^2/(2t)+(k-4)\log2}}
                {(k-4)/(2t)-\log2}\\
 &\le\frac{2t}{k-6}e^{-(k-4)^2+(k-4)\log2}.
\end{align*}
Both denominators are positive, and the last step follows from
$2t\le1$, $2t\log2\le2$. Therefore
\[
 \delta_3\le\frac{\sqrt t}{\sqrt\pi}
 \sum_{k\ge b}\frac{f_k}{(k-6)a^k},\qquad
 f_k=(1.1)^k\Gamma(k/2)e^{-(k-4)^2+(k-4)\log2}.
\]
The two bases are
\[
 f_{14}=(1.1)^{14}6!2^{10}e^{-100}<1.042\cdot10^{-37},
\]
\[
 f_{15}=(1.1)^{15}\frac{14!\sqrt\pi}{4^7 7!}2^{11}e^{-121}
                                         <4.516\cdot10^{-46}.
\]
Exact gamma recurrence gives
\[
 \frac{f_{k+2}}{f_k}=2.42k e^{12-4k}<1\qquad(k\ge14).
\]
Indeed $x e^{-4x}$ decreases for $x\ge14$ and the ratio at $14$
is $33.88e^{-44}<1$, already implied by $e^{44}>45$.
Parity induction proves $f_k<10^{-30}$ for every integer $k\ge14$.
Since $\sqrt t/\sqrt\pi<1$ and $1/(k-6)<1$ throughout the
original range, enlarging the nonnegative sum now gives
\begin{equation}\label{tailfinal}
 \delta_3\le\sum_{k\ge14}\frac{10^{-30}}{a^k}
 =\frac{10^{-30}}{a^{14}(1-a^{-1})}
 \le\frac{2\cdot10^{-30}}{a^{14}}.
\end{equation}
This is an inequality, with no altered identity or cutoff.

For $0\le\sigma\le1$ and $a>3$ the available denominator slack is
\begin{align*}
 \frac{0.031\,2^{-\sigma}}{a-0.865}
 -\frac{0.031\,2^{-\sigma}}{a-0.353}
 &=\frac{0.031(0.512)2^{-\sigma}}
          {(a-0.865)(a-0.353)}\\
 &\ge\frac{0.007936}{a^2}
 >\frac{2\cdot10^{-30}}{a^{14}}.
\end{align*}
The last step follows from $a^{12}>3^{12}$ and
$0.007936\,3^{12}>2\cdot10^{-30}$. Multiplying the slack by
$e^{3.49/(T-4)}>1$ only increases it. Equations \eqref{deltaone},
\eqref{deltatwo}, and \eqref{tailfinal} consequently give
\[
 \delta\le\frac{0.366\,9^\sigma+0.031\,2^{-\sigma}}{a-0.865}
                       e^{3.49/(T-4)}
 \le\frac{0.397\,9^\sigma}{a-0.865}e^{3.49/(T-4)}.
\]
Together with \eqref{epsdelta} and \eqref{epsilonfinal}, this proves
Theorem \ref{main} for $t>0$.

\section{The endpoint $t=0$}

At $t=0$ use $T'=T$, $u=\sigma$, and $K=1$ directly in
\eqref{RSexp}. The pointwise estimate \eqref{Ebound} and the
nonnegative-$u$ estimate above give
\[
 \left|\frac{R_{0,N}(\sigma+iT)}
 {(-1)^{N-1}Ue^{\pi i\sigma/4}M_0(iT)}-C_0(p)\right|
 \le\tfrac12(e^h-1)+e^h\frac{0.2\,9^\sigma}{a-0.865},
 \quad h=\frac{\sigma^2+c}{T-3}.
\]
Here
\[
 0\le h\le\frac{23}{12(T-3)}
 \le\frac{10}{3(T-6)}\le\frac{3.49}{T-4}.
\]
The first contribution is at most
$5e^{3.49/(T-4)}/(3(T-6))$ by $e^h-1\le he^h$;
the second is at most
$0.397\,9^\sigma e^{3.49/(T-4)}/(a-0.865)$ because $0.2<0.397$.
This proves the endpoint without passing a limit through $\lfloor a\rfloor$.

\section{Exact rational certificates for the finite comparisons}

The companion \texttt{exact\_certificates.py} uses Python integer
arithmetic, \texttt{Fraction}, and factorials only. Every comparison
is exact. No floating-point comparison is used. Its rational
enclosures follow from these elementary series arguments.

For $0<z<1$, alternating integration of the finite geometric series
gives
$\arctan z=\sum_{j\ge0}(-1)^jz^{2j+1}/(2j+1)$,
between a partial sum and that sum plus its next term. Machin's
identity $\pi=16\arctan(1/5)-4\arctan(1/239)$ follows from
the tangent addition formula and the fact that the relevant angle
is in $(0,\pi/2)$: $\tan(2\arctan(1/5))=5/12$,
$\tan(4\arctan(1/5))=120/119$, and subtracting
$\arctan(1/239)$ gives tangent $1$. Using partial sums through
$j=40$ and $j=10$, respectively, proves
\[
 3.1415926535<\pi<3.1415926536,
 \qquad \sqrt\pi<1.772453851,
\]
where the square-root assertion is checked by squaring positive
rationals.

For $x>1$, $z=(x-1)/(x+1)$ satisfies
\[
 \log x=2\sum_{j=0}^n\frac{z^{2j+1}}{2j+1}+R_n,
 \qquad 0<R_n\le\frac{2z^{2n+3}}{(2n+3)(1-z^2)}.
\]
This follows by integrating the geometric expansion of
$2/(1-z^2)$ and bounding every remaining denominator below by
$2n+3$. With $n=60$ for $x=2,3$, these give the stated enclosures
for logarithms and hence for $A$.

For rational $x\ge0$ and integer $n$ with $x<n+2$, write
$E_n(x)=\sum_{j=0}^n x^j/j!$. The ratios of successive terms
after the first omitted term are at most $x/(n+2)$, whence
\[
 E_n(x)\le e^x\le E_n(x)+
 \frac{x^{n+1}}{(n+1)!}\frac1{1-x/(n+2)}.
\]
The script uses $n=40$ for the two small positive exponential
arguments, and lower sums with $n=300$ at $x=100$ and $n=330$
at $x=121$ for the negative-exponential tail bases. This proves
all comparisons used above, in particular
\[
 3.5670161955<A<3.5670161956,
\]
\[
 0.2\exp\!\left((\log3)^2/2+(A-3.49)/96\right)<0.365986,
 \qquad 0.029e^{(\log\sqrt2)^2/2}<0.031.
\]
The script's 23 explicit checks all passed in both ordinary and
optimized Python mode on 20 September 2026. Failed comparisons and
invalid series domains raise exceptions explicitly, so optimization
cannot disable the checks.
They supplement the recurrence and monotonicity proofs; they do
not substitute a finite sample for a bound over the unbounded
integer indices or continuous parameter ranges.

\begin{thebibliography}{9}
\bibitem{polymath} D.~H.~J. Polymath,
\emph{Effective approximation of heat flow evolution of the Riemann
$\xi$ function, and a new upper bound for the de Bruijn--Newman
constant}, \href{https://arxiv.org/abs/1904.12438v2}{arXiv:1904.12438v2}.
Original author source \texttt{debruijn.tex}; exact labels and source
coverage are specified in the text.
\bibitem{arias} J. Arias de Reyna,
\emph{High-precision computation of Riemann's zeta function by the
Riemann--Siegel asymptotic formula, I}, Mathematics of Computation
\textbf{80} (2011), 995--1009. Cited here through the precise input
statement and attribution in Polymath's original source.
\bibitem{boyd} W.~G.~C. Boyd,
\emph{Gamma Function Asymptotics by an Extension of the Method of
Steepest Descents}, Proceedings: Mathematical and Physical Sciences
\textbf{447}, no.~1931 (1994), 609--630. Cited here through
Polymath's source Lemma \texttt{elem-lem}(v).
\end{thebibliography}
\end{document}
